\BeleskeID{09_2-s017}
% element: 09_2-s017:e04
N tranzistor tek počinje da vodi malu struju, a visok izlazni napon obezbeđuje zasićenje. Ravnoteža struja u oznakama priključaka glasi
\[\frac{k_p}2\left(2V_{DS,P}(V_{GS,P}-V_{Tp})-V_{DS,P}^2\right)=\frac{k_n}2(V_{GS,N}-V_{Tn})^2\]
Zamenjuju se $V_{DS,P}=V_O-V_{DD}$, $V_{GS,P}=-V_{DD}$ i $V_{GS,N}=V_I$. Pri određivanju donjeg ulaznog praga diferencira se ravnoteža struja po ulaznom naponu, a zatim zamenjuje nagib $-1$:
\[k_p\left(-2(V_{DD}+V_{Tp})-(V_O-V_{DD})\right)-k_p(V_O-V_{DD})=2k_n(V_I-V_{Tn})\]
\[k_p\left(-2(V_{DD}+V_{Tp})-2(V_O-V_{DD})\right)=2k_n(V_I-V_{Tn1})\]
Sistem čine ravnoteža struja i jednačina dobijena iz nagiba. Njihovo zajedničko rešavanje daje ulazni i izlazni napon u istoj tački karakteristike.
